\documentclass[11pt,a4paper]{article}
\usepackage[margin=1in]{geometry}
\usepackage{hyperref}
\usepackage{enumitem}
\usepackage{graphicx}
\usepackage{xcolor}

% For tables
\usepackage{booktabs}

% Variable: Lab Instructor
\makeatletter \newcommand{\BvrSetLabInstructor}[1]{%
  \gdef\Bvr@LabInstructor{#1}}
% default
\newcommand{\Bvr@LabInstructor}{Ms Paramveer %
  Sidhu}
% public accessor
\newcommand{\BvrLabInstructorValue}{\Bvr@LabInstructor}
\makeatother

% Add subtitle
\usepackage{titling}
% -- Set title bold
\pretitle{\begin{center}\bfseries\fontsize{18bp}{18bp}\selectfont}
\makeatletter
\newcommand{\BvrSubtitle}[1]{%
  \posttitle{%
    \par\end{center}
    \begin{center}\large#1\end{center}
    \vskip 0.5em}%
}
\makeatother

% Hints in the section title(s)
\newcommand{\BvrHint}[1]{%
  {\normalfont\normalsize\itshape\hfill #1}}

% Table formatting
\newcommand{\BvrTableHeader}[1]{\textbf{#1}}

% =====================================================

\title{Post-Hospital Recovery Management and Early-Warning System}

\BvrSetLabInstructor{Ms Paramveer Sidhu}

\BvrSubtitle{%
  Project Proposal for UCS503P  \\
  Submitted to: \BvrLabInstructorValue \\ \bfseries
  Thapar Institute of Engineering and Technology% 
}

\author{
  Pulkit Sareen, CSED \and
  Shivam Singh, CSED \and
  Daksh Anand, CSED%
}

\date{16 August, 2026}

\begin{document}
\maketitle

\section{Higher-order goal%
  \BvrHint{\ldots secondary framing}}

This project aims to improve the continuity and organization
of patient care after hospital discharge by providing a centralized software
platform for monitoring recovery progress, managing post-discharge activities,
and identifying unusual recovery patterns that may require professional
attention.

The immediate engineering objective is to translate
post-hospital recovery management into a measurable, deployable software
solution that connects patients, healthcare professionals, and recovery data
through a structured workflow.

The system will not attempt to diagnose diseases or replace
healthcare professionals. Instead, it will provide decision-support features
that organize patient-reported information, monitor recovery trends, and
highlight potentially concerning deviations for professional review.

The project combines conventional software engineering with a
lightweight machine-learning component. The software platform will manage
users, recovery plans, medications, exercises, appointments, daily logs,
notifications, and historical records, while machine learning will be used to
identify deviations from expected recovery patterns and classify recovery
trajectories.

\section{Time-to-value %
\BvrHint{\ldots primary heuristic}}

We will deliver the project through incremental iterations,
beginning with the core patient recovery workflow before introducing
machine-learning functionality.

The first iteration will focus on establishing a usable
system where:
\begin{itemize}[leftmargin=0.7cm]
    \item A healthcare professional can create a patient's post-discharge recovery plan.
    \item A patient can access the plan and submit daily recovery information.
    \item Medication, exercise, and appointment schedules can be tracked.
    \item Recovery information is stored with timestamps.
    \item Patients can view their recovery history and progress.
    \item Healthcare professionals can view patient recovery information.
\end{itemize}

Subsequent iterations will introduce:
\begin{itemize}[leftmargin=0.7cm]
    \item Recovery scoring.
    \item Recovery trend visualization.
    \item Rule-based alerts.
    \item Machine-learning-based trajectory analysis.
    \item Risk/deviation indicators.
    \item Healthcare-professional dashboards.
    \item Notifications.
\end{itemize}

We will prioritize working software and measurable
functionality over premature complexity. Weekly iterations will be used to
obtain feedback and identify problems in the workflow early.

Success will be measured not only by whether the ML model
produces predictions, but also by whether the complete software system can
reliably collect recovery information, identify changes, present them clearly,
and support timely professional review.

\section{Problem Statement}

After a patient leaves a hospital, recovery continues for
days or weeks, but the patient's recovery information is often fragmented
between prescriptions, appointments, personal notes, and occasional follow-up
consultations.

Common problems include:
\begin{itemize}[leftmargin=0.7cm]
    \item \textbf{Limited continuous monitoring:} Healthcare professionals may have little visibility into the patient's recovery between scheduled appointments.
    \item \textbf{Fragmented information:} Medication schedules, exercise plans, symptoms, pain levels, mobility, and appointments may be tracked separately.
    \item \textbf{Poor visibility of recovery trends:} A single measurement may not indicate whether a patient is improving or deteriorating. Changes over several days can be more informative.
    \item \textbf{Delayed identification of unusual patterns:} A patient may experience a gradual deviation from their expected recovery pattern without recognizing its significance.
    \item \textbf{Medication and exercise adherence:} Patients may forget or inconsistently follow prescribed post-discharge activities.
    \item \textbf{Follow-up coordination:} Appointments and recovery milestones can be difficult to manage alongside daily recovery tasks.
    \item \textbf{Information overload for healthcare professionals:} If large amounts of patient data are collected, professionals need a way to prioritize patients requiring attention rather than manually inspecting every record.
    \item \textbf{Lack of structured historical data:} Recovery information collected from patients is often not organized into a consistent longitudinal record.
\end{itemize}

Why this matters now: post-hospital care involves a period
where the patient is responsible for a significant part of their recovery while
professional interaction may be intermittent. A software system that organizes
this period into a continuous and measurable workflow can reduce information
gaps and help professionals identify patients whose recovery requires review.

\section{Proposed Solution %
\BvrHint{\ldots Software Proposition}}

\subsection{Overview}

Build a web-based Post-Hospital Recovery Management and
Early-Warning System with the following components:
\begin{itemize}[leftmargin=0.7cm]
    \item \textbf{Patient portal:} Patients view their recovery plan, medications, exercises, appointments, and daily tasks.
    \item \textbf{Daily recovery logging:} Patients record structured information such as pain, mobility/function, sleep, symptoms, medication adherence, and exercise completion.
    \item \textbf{Recovery timeline:} The system displays the patient's recovery progress over time instead of only showing the latest measurement.
    \item \textbf{Recovery scoring:} A transparent application-level score summarizes multiple recovery dimensions.
    \item \textbf{Rule-based monitoring:} Explicitly configured thresholds and repeated changes generate alerts for review.
    \item \textbf{ML-based trajectory analysis:} A machine-learning model analyzes combinations and trends in recovery data to identify patterns that deviate from expected recovery trajectories.
    \item \textbf{Healthcare-professional dashboard:} Professionals can view patient status, trends, alerts, and historical recovery information.
    \item \textbf{Medication and exercise management:} Recovery activities can be scheduled and adherence can be recorded.
    \item \textbf{Appointment management:} Follow-up appointments and recovery milestones can be stored and tracked.
    \item \textbf{Notification system:} Patients can receive reminders while healthcare professionals can receive high-priority review notifications.
    \item \textbf{Administrative dashboard:} Provides system-level information such as active patients, outstanding alerts, adherence statistics, and recovery trends.
\end{itemize}

The system is designed as a decision-support platform rather
than an autonomous diagnostic system. ML predictions and alerts will indicate
that a patient's data requires review rather than claiming that a patient has a
particular disease or complication.

\subsection{Core workflow}

\subsubsection{1. Recovery plan creation}

A healthcare professional creates a recovery plan for a
discharged patient.

The plan contains:
\begin{itemize}[leftmargin=0.7cm]
    \item Recovery category/procedure
    \item Expected recovery duration
    \item Recovery milestones
    \item Measurements to track
    \item Medication schedule
    \item Exercise/activity schedule
    \item Restrictions
    \item Follow-up appointments
\end{itemize}

\subsubsection{2. Patient onboarding}

The patient receives access to the recovery plan and can view
their daily tasks.

\subsubsection{3. Daily check-in}

The patient submits structured information such as:
\begin{itemize}[leftmargin=0.7cm]
    \item Pain level
    \item Mobility/function level
    \item Sleep duration
    \item Temperature, when relevant to the selected recovery plan
    \item Medication adherence
    \item Exercise adherence
    \item Newly reported symptoms
    \item General recovery notes
\end{itemize}

The system records every submission with a timestamp.

\subsubsection{4. Recovery analysis}

The system calculates the patient's current recovery
indicators and compares recent measurements with historical measurements and
the expected recovery trajectory.

\subsubsection{5. Rule-based monitoring}

Explicitly configured conditions are checked.

For example:
\begin{itemize}[leftmargin=0.7cm]
    \item A measurement exceeds a configured threshold.
    \item A concerning symptom is reported.
    \item Medication is repeatedly missed.
    \item Recovery measurements change significantly over a short period.
\end{itemize}

These events can create a review flag.

\subsubsection{6. Machine-learning analysis}

The ML model analyzes the patient's longitudinal data and
estimates whether the current trajectory is:
\begin{itemize}[leftmargin=0.7cm]
    \item On Track
    \item Delayed
    \item Potentially Concerning
\end{itemize}

The model may also produce a deviation/risk probability.

\subsubsection{7. Professional review}

Healthcare professionals see prioritized patients on their
dashboard.

They can inspect:
\begin{itemize}[leftmargin=0.7cm]
    \item Current values
    \item Historical trends
    \item Recovery score
    \item ML output
    \item Reasons/features associated with the flag
    \item Medication adherence
    \item Exercise adherence
    \item Patient notes
    \item Previous alerts
\end{itemize}

The healthcare professional remains responsible for deciding
whether further action is necessary.

\subsubsection{8. Feedback}

The outcome of professional review can be recorded in the
system.

This information can later become labeled data for evaluating
or improving future versions of the ML model.

\section{Solution Approach %
\BvrHint{\ldots Engineering focus}}

This is an engineering course project, so the primary focus
will be on building a reliable software system around the recovery workflow
rather than developing a novel medical ML algorithm.

The ML component will use established, interpretable
machine-learning techniques on structured longitudinal data.

\subsection{Recovery data model %
\BvrHint{\ldots structured patient history}}

The system will represent recovery as a sequence of
time-stamped observations rather than isolated values.

A daily recovery record may contain:

\begin{center}
\begin{tabular}{ll}
\toprule
\textbf{Feature} & \textbf{Example} \\
\midrule
Recovery day & Day 7 \\
Pain & 5/10 \\
Mobility/function & 62\% \\
Sleep & 6.5 hours \\
Medication adherence & 100\% \\
Exercise adherence & 67\% \\
Temperature & 37.2°C \\
Reported symptoms & None \\
Patient note & Optional \\
Timestamp & Date/time \\
\bottomrule
\end{tabular}
\end{center}

The exact measurements will depend on the selected recovery
scenario.

The initial implementation will focus on a limited set of
measurable variables rather than attempting to monitor every possible medical
parameter.

\subsection{Recovery scoring %
\BvrHint{\ldots transparent logic}}

A preliminary recovery score will be implemented using a
transparent weighted formula.

For example:

Recovery Score =
\begin{itemize}[leftmargin=0.7cm]
    \item 30\% Mobility/Function
    \item 25\% Pain-related indicator
    \item 15\% Exercise adherence
    \item 15\% Medication adherence
    \item 15\% Sleep indicator
\end{itemize}

The weights are application-level prototype assumptions and
will not be presented as medically validated clinical standards.

The purpose of this score is to provide an understandable
summary of multiple recovery dimensions.

It will be complemented by the patient's individual trend
rather than being used as a standalone medical decision.

\subsection{Rule-based alerts %
\BvrHint{\ldots deterministic safety layer}}

Rules will be used for conditions that can be explicitly
specified.

For example:
\begin{itemize}[leftmargin=0.7cm]
    \item \textbf{IF} a configured critical symptom is reported $\rightarrow$ create high-priority review flag
    \item \textbf{IF} a monitored measurement exceeds its configured threshold $\rightarrow$ create monitoring alert
    \item \textbf{IF} medication adherence falls below a configured level $\rightarrow$ create adherence alert
    \item \textbf{IF} recovery measurements show repeated deterioration $\rightarrow$ increase monitoring priority
\end{itemize}

The rules will be configurable according to the selected
recovery scenario.

The system will not autonomously diagnose the patient or
prescribe treatment.

\subsection{Machine-learning component %
\BvrHint{\ldots non-research approach}}

Machine learning will focus on recovery trajectory
classification and anomaly detection.

The initial model will use structured features such as:
\begin{itemize}[leftmargin=0.7cm]
    \item Current pain level
    \item Change in pain over previous days
    \item Mobility/function
    \item Change in mobility
    \item Sleep
    \item Medication adherence
    \item Exercise adherence
    \item Reported symptoms
    \item Recovery day
    \item Procedure/recovery category
\end{itemize}

The initial classification task will be:
\begin{itemize}[leftmargin=0.7cm]
    \item ON TRACK
    \item DELAYED
    \item POTENTIALLY CONCERNING
\end{itemize}

A lightweight model such as Random Forest will be
evaluated first because it works well with tabular data and is relatively
interpretable.

For anomaly detection, an approach such as Isolation
Forest may be evaluated to identify recovery observations that differ
substantially from learned patterns.

The project will compare the performance of the selected
model against a simple rule-based baseline.

This allows the project to answer an important engineering
question: Does adding machine learning provide useful information
beyond a simple deterministic recovery-score/rule system?

\subsection{Dataset strategy %
\BvrHint{\ldots practical educational approach}}

The project will not depend on finding one perfect public
post-discharge dataset.

The dataset strategy will use two sources where possible:

\subsubsection{Public/real-world data}

Relevant publicly available healthcare and
patient-reported-outcome datasets will be investigated and used where their
variables and licensing/access conditions are appropriate.

\subsubsection{Synthetic longitudinal data}

Because a clean dataset containing all required
post-discharge variables may not be publicly available, a synthetic patient
dataset will be generated for development and testing.

The synthetic data will represent different recovery
trajectories, such as:
\begin{itemize}[leftmargin=0.7cm]
    \item Normal recovery
    \item Gradually improving recovery
    \item Delayed recovery
    \item Recovery with adherence problems
    \item Recovery with temporary deviations
    \item Recovery with significant deterioration patterns
\end{itemize}

The synthetic data generator will create time-series
observations rather than independent random records.

This allows the system to contain realistic longitudinal
relationships such as:

Expected recovery: Pain: $8 \rightarrow 7 \rightarrow 6 \rightarrow 5 \rightarrow 4 \rightarrow 3$

Mobility: $30 \rightarrow 38 \rightarrow 45 \rightarrow 55 \rightarrow 65 \rightarrow 72$

while also generating patients whose trajectories deviate
from the expected pattern.

Synthetic data will be clearly identified as simulated data
and will not be presented as real clinical evidence.

\subsection{Web architecture %
\BvrHint{\ldots web-based solution}}

The system will use a 3-tier architecture:

\subsubsection{Frontend}

A responsive React application will provide:
\begin{itemize}[leftmargin=0.7cm]
    \item Patient dashboard
    \item Daily check-in
    \item Recovery timeline
    \item Medication tracker
    \item Exercise tracker
    \item Appointment view
    \item Healthcare-professional dashboard
    \item Alert management
\end{itemize}

\subsubsection{Backend}

A FastAPI backend will handle:
\begin{itemize}[leftmargin=0.7cm]
    \item Authentication
    \item Authorization
    \item Patient management
    \item Recovery plans
    \item Daily logs
    \item Medication schedules
    \item Exercise schedules
    \item Appointment management
    \item Alert generation
    \item ML prediction requests
    \item Notification logic
    \item Audit/history records
\end{itemize}

\subsubsection{Database}

The database will store:
\begin{itemize}[leftmargin=0.7cm]
    \item Users
    \item Patients
    \item Healthcare professionals
    \item Recovery plans
    \item Daily recovery logs
    \item Medications
    \item Exercises
    \item Appointments
    \item Alerts
    \item ML predictions
    \item Review outcomes
\end{itemize}

The architecture will keep the ML component separate from the
primary API logic so that the model can be updated without redesigning the
entire application.

\subsection{ML service architecture}

The ML pipeline will be separated into:

Patient Recovery Data $\rightarrow$
Data Validation $\rightarrow$
Feature Engineering $\rightarrow$
ML Model $\rightarrow$
Prediction $\rightarrow$
Recovery Engine $\rightarrow$
Database $\rightarrow$
Dashboard / Alert

The model will not directly communicate a medical diagnosis
to the patient.

Instead, it will return structured information such as:
\begin{itemize}[leftmargin=0.7cm]
    \item status: ``potentially\_concerning''
    \item probability: 0.76
    \item model\_version: ``v1''
    \item timestamp: ...
\end{itemize}

The backend can then determine how that prediction should be
displayed and whether professional review should be prioritized.

\subsection{Authentication and authorization}

Different user roles will have different permissions.

\subsubsection{Patient}

Can:
\begin{itemize}[leftmargin=0.7cm]
    \item view their recovery plan
    \item submit daily logs
    \item view their own recovery history
    \item manage reminders
    \item view appointments
    \item receive notifications
\end{itemize}

\subsubsection{Healthcare professional}

Can:
\begin{itemize}[leftmargin=0.7cm]
    \item create recovery plans
    \item view assigned patients
    \item inspect recovery trends
    \item review alerts
    \item record review outcomes
    \item modify recovery schedules
\end{itemize}

\subsubsection{Administrator}

Can:
\begin{itemize}[leftmargin=0.7cm]
    \item manage users
    \item monitor system activity
    \item view aggregate statistics
    \item manage system configuration
\end{itemize}

Role-based access control will prevent patients from
accessing other patients' information.

\section{Evaluation Criterion %
\BvrHint{\ldots measurable and attributable}}

The project will define measurable criteria for both the
software system and the ML component.

\subsection{Primary evaluation metric}

\textbf{Recovery monitoring completeness}

Measure the percentage of scheduled patient check-ins that
are successfully recorded and processed by the system.

\textbf{Target:}
$\ge$ 90\% of simulated scheduled check-ins successfully stored
and processed during the pilot.

\textbf{Attribution:}
Calculated from database records containing expected
check-ins and successfully submitted logs.

\subsection{Secondary software metrics}

\textbf{Alert processing time}

Measure the time between submission of a recovery log and
creation of the corresponding alert/prediction.

\textbf{Target:} The system should process a normal recovery submission within
a few seconds under the project test load.

\textbf{Medication adherence tracking}

Measure whether the system correctly records scheduled and
completed medication events.

\textbf{Appointment tracking}

Measure whether scheduled appointments and reminders are
correctly generated and displayed.

\textbf{System reliability}

Measure:
\begin{itemize}[leftmargin=0.7cm]
    \item API availability
    \item failed requests
    \item database errors
    \item authentication failures
\end{itemize}

\textbf{User clarity}

Pilot users can rate how clearly the system communicates:
\begin{itemize}[leftmargin=0.7cm]
    \item recovery status
    \item daily tasks
    \item alerts
    \item recovery trends
\end{itemize}

using a simple 1--5 scale.

\subsection{ML evaluation metrics}

The ML model will be evaluated using:
\begin{itemize}[leftmargin=0.7cm]
    \item Accuracy
    \item Precision
    \item Recall
    \item F1-score
    \item Confusion matrix
\end{itemize}

Special attention will be given to recall for the
potentially concerning class, because failing to identify a concerning
simulated recovery trajectory is an important failure mode for the prototype.

The model will be compared with a simple rule-based baseline.

Example evaluation:

\begin{center}
\begin{tabular}{lcc}
\toprule
& \textbf{Rule System} & \textbf{ML Model} \\
\midrule
Accuracy & XX\% & XX\% \\
Precision & XX\% & XX\% \\
Recall & XX\% & XX\% \\
F1 Score & XX\% & XX\% \\
\bottomrule
\end{tabular}
\end{center}

The final values will be filled after model training and
testing rather than being claimed beforehand.

\subsection{Pilot validation plan}

A controlled educational pilot can be conducted using:
\begin{itemize}[leftmargin=0.7cm]
    \item 10--20 test users
    \item Simulated patient recovery cases
    \item Student testers acting as patients
    \item Student/teacher-supervised users acting as healthcare professionals where appropriate
\end{itemize}

The pilot will measure:
\begin{itemize}[leftmargin=0.7cm]
    \item Completion of daily check-ins
    \item Correctness of workflow
    \item Alert generation
    \item Dashboard usability
    \item ML prediction performance
    \item Response time
    \item User clarity
\end{itemize}

No real patient's medical information will be required for
the educational prototype.

\section{Scalability %
\BvrHint{\ldots with foresight}}

The proposition should scale in both theoretical and
practical senses.

\subsection{1. Scalable architecture}

The system will support growth in patient and recovery-log
volume through:
\begin{itemize}[leftmargin=0.7cm]
    \item Separate frontend and backend components
    \item Stateless API design
    \item Database indexing
    \item Pagination for patient records
    \item Efficient querying of time-series recovery logs
    \item Separate ML inference service
    \item Background processing for notifications
\end{itemize}

\subsection{2. ML scalability}

The ML model will be versioned independently from the main
application.

For example:

Model v1 $\rightarrow$ Model v2 $\rightarrow$ Model v3

Predictions will store the model version that generated them.

This makes it possible to compare model versions and roll
back an updated model if required.

\subsection{3. Operational deployment}

The system should be deployable using common web
infrastructure:
\begin{itemize}[leftmargin=0.7cm]
    \item Managed database
    \item Cloud-hosted backend
    \item Cloud-hosted frontend
    \item Containerized services where appropriate
    \item CI/CD pipeline
    \item Environment-based configuration
\end{itemize}

\section{Engine availability heuristic}

A mature set of development tools and services is available
to implement the system as a web-based project.

The proposed stack is:
\begin{itemize}[leftmargin=0.7cm]
    \item \textbf{Frontend:} React
    \item \textbf{Styling/UI:} Tailwind CSS
    \item \textbf{Backend:} FastAPI
    \item \textbf{Database:} PostgreSQL or MongoDB
    \item \textbf{ML:} Python, NumPy, pandas, scikit-learn
    \item \textbf{Data visualization:} Recharts
    \item \textbf{Authentication:} JWT/token-based authentication
    \item \textbf{Testing:} Pytest and frontend testing tools
    \item \textbf{Version control:} Git and GitHub
    \item \textbf{CI:} GitHub Actions
    \item \textbf{Deployment:} Vercel for frontend and Render or equivalent platform for backend
    \item \textbf{Containerization:} Docker where required
\end{itemize}

These are mature technologies, allowing the project team to
focus on recovery workflow design, ML integration, testing, security, and
system correctness rather than developing infrastructure from scratch.

\section{Project scope and deliverables %
\BvrHint{\ldots iteration-friendly}}

\subsection{Initial deliverable %
\BvrHint{\ldots first iteration}}

The first iteration will contain:
\begin{itemize}[leftmargin=0.7cm]
    \item Patient and healthcare-professional authentication
    \item Patient profile
    \item Recovery plan creation
    \item Patient recovery-plan view
    \item Daily recovery check-in
    \item Recovery data storage
    \item Medication schedule
    \item Exercise schedule
    \item Appointment schedule
    \item Basic recovery history
    \item Database timestamps
    \item Basic backend tests
    \item CI pipeline
    \item Deployment to staging
\end{itemize}

At the end of this iteration, a complete patient $\rightarrow$ daily
check-in $\rightarrow$ database $\rightarrow$ professional dashboard workflow should be functional.

\subsection{Subsequent deliverables}

\subsubsection{Iteration 2}
\begin{itemize}[leftmargin=0.7cm]
    \item Recovery score
    \item Recovery trend graphs
    \item Medication adherence calculation
    \item Exercise adherence calculation
    \item Rule-based alerts
    \item Notification system
\end{itemize}

\subsubsection{Iteration 3}
\begin{itemize}[leftmargin=0.7cm]
    \item Synthetic patient-data generator
    \item ML training pipeline
    \item Recovery trajectory classification
    \item Anomaly detection
    \item Model evaluation
    \item Prediction API
\end{itemize}

\subsubsection{Iteration 4}
\begin{itemize}[leftmargin=0.7cm]
    \item Healthcare-professional command center
    \item Patient prioritization
    \item ML prediction visualization
    \item Alert history
    \item Professional review workflow
    \item Feedback/outcome recording
\end{itemize}

\subsubsection{Final iteration}
\begin{itemize}[leftmargin=0.7cm]
    \item Security improvements
    \item Performance testing
    \item Integration testing
    \item ML model comparison
    \item UI refinement
    \item Documentation
    \item CI/CD improvements
    \item Final deployment
    \item Final evaluation
\end{itemize}

\section{Risks and mitigations}

\subsection{Medical interpretation risk}

\textbf{Risk:} Users may interpret the ML output as a medical
diagnosis.

\textbf{Mitigation:}
The system will clearly state that ML outputs are
decision-support indicators for professional review. The system will not
provide autonomous diagnosis or treatment recommendations.

\subsection{Data availability}

\textbf{Risk:} A public dataset may not contain all required
post-discharge variables.

\textbf{Mitigation:}
Use relevant public datasets where appropriate and create a
documented synthetic longitudinal dataset for development and evaluation.

\subsection{Synthetic-data bias}

\textbf{Risk:} A model trained primarily on synthetic data may not
represent real-world patients.

\textbf{Mitigation:}
Clearly identify synthetic data as simulated, evaluate the
model only within the scope of the simulated experiment, and avoid making
claims of clinical validity.

\subsection{False alerts}

\textbf{Risk:} The system may produce too many alerts.

\textbf{Mitigation:}
Use a combination of deterministic rules and ML confidence
thresholds, record alert history, and evaluate false-positive rates during
testing.

\subsection{Missed concerning patterns}

\textbf{Risk:} The ML system may fail to identify a simulated
concerning recovery pattern.

\textbf{Mitigation:}
Evaluate recall separately for the potentially concerning
class and keep the professional review layer in the workflow.

\subsection{Data privacy}

\textbf{Risk:} Recovery information is sensitive.

\textbf{Mitigation:}
For the educational prototype, use synthetic data whenever
possible. If real data is ever introduced, collect only necessary information,
apply role-based access control, avoid exposing patient information through
logs, and follow appropriate privacy requirements.

\subsection{Model drift}

\textbf{Risk:} A future dataset may differ from the training data.

\textbf{Mitigation:}
Store model versions, prediction timestamps, and evaluation
metrics. Future models will be evaluated before deployment.

\subsection{System reliability}

\textbf{Risk:} Failure of the ML service should not make the entire
application unavailable.

\textbf{Mitigation:}
The core patient-recording and recovery-management system
will continue to function even if ML inference is temporarily unavailable. ML
predictions will be treated as an additional service rather than the foundation
of the entire application.

\section{Summary}

This project proposes a deployable web-based Post-Hospital
Recovery Management and Early-Warning System that provides a structured
workflow for managing the period after hospital discharge.

The system will allow healthcare professionals to create
recovery plans and patients to record daily recovery information, manage
medications and exercises, track appointments, and monitor their progress over
time.

A transparent recovery-score and rule-based monitoring layer
will provide understandable baseline functionality, while machine learning will
analyze longitudinal recovery data to identify patterns that deviate from
expected recovery trajectories.

The project deliberately avoids attempting autonomous medical
diagnosis. Instead, ML outputs will be presented as indicators that can help
prioritize cases for professional review.

The project meets the requirements of a Software Engineering
course by emphasizing:
\begin{itemize}[leftmargin=0.7cm]
    \item Requirements and workflow design
    \item Role-based access control
    \item Web architecture
    \item Database design
    \item API development
    \item Testing
    \item CI/CD
    \item Deployment
    \item Scalability
    \item Measurable evaluation
    \item Machine-learning integration
    \item Data privacy
    \item Human-in-the-loop decision making
\end{itemize}

The project will follow an iterative development process,
beginning with the core recovery-management workflow and progressively adding
analytics, alerts, machine learning, and professional dashboards.

The final system will demonstrate how a relatively
lightweight machine-learning component can be integrated into a complete
software product to transform longitudinal recovery data into useful,
explainable, and actionable information for professional review.

\end{document}
